\section{Verification purpose}

This case is a negative verification test for mesh dependence in the current
multiple-discrete-cell treatment of structures. A physical
\SI{10}{m}-by-\SI{10}{m} house may occupy one raster cell, several independent
cells, one oversized cell, or disappear under categorical resampling as the
mesh changes. Because those cells are not assembled into one conserved
structure object, changing resolution alters represented heat release,
firebrand generation, deposition, and structure-to-structure ignition.

The test asks whether all 15 input-matched simulations produce an
evaluable community rate of spread (ROS), and whether the fitted values remain
within the predeclared mesh-independence limits relative to the run-derived
10-m anchor. Detecting the known resolution sensitivity is useful diagnostic
evidence, but it is a \emph{verification failure}, not a pass.

\section{Mathematical formulation and numerical implementation}

\subsection{Mode-resampled square array}

The physical layout is a periodic array of square structures of side
\(a=\SI{10}{m}\), separated in both coordinates by
\(d=\SI{10}{m}\). Define the one-dimensional structural indicator
\[
 I(s)=\begin{cases}1,&s\bmod(a+d)<a,\\0,&\text{otherwise.}\end{cases}
\]
For grid interval \(C_i=[i\Delta,(i+1)\Delta]\), preprocessing calculates
\[
 f_i=\frac{1}{\Delta}\int_{C_i}I(s)\,ds,
 \qquad b_i=\mathbf{1}(f_i\ge 1/2),
 \qquad B_{ij}=b_i b_j .
\]
This coordinate-wise categorical-mode construction reproduces the reported
geometries. At \(\Delta=\SI{7}{m}\), effective side lengths are 7 or
\SI{14}{m} and footprints range from 49 to \(\SI{196}{m^2}\). At
\(\Delta=\SI{10}{m}\), footprint and gap are exact. At
\(\Delta=\SI{30}{m}\), one represented structure cell has area
\SI{900}{m^2} and represented structures are separated by \SI{30}{m}.
The 20- and 40-m maps eliminate represented firebreaks; this is an intended
resampling diagnostic, not a preprocessing failure.

\subsection{Cell-local firebrand treatment}

A represented structural cell uses a design-fire curve \(q^{\prime\prime}(t)\). The curve has a 300-s linear growth phase, a one-hour plateau at \SI{400}{kW.m^{-2}}, and a 300-s decay. Under the current
multiple-discrete-cell treatment,
\[
 \mathrm{HRR}_{ij}(t)=q^{\prime\prime}(t)\Delta^2,
 \qquad \dot N_{ij}(t)=GR\,\frac{\mathrm{HRR}_{ij}(t)}{1000},
 \qquad GR=\SI{10}{pcs.MW^{-1}.s^{-1}}.
\]
Changing \(\Delta\) therefore changes both represented geometry and the heat-release rate (HRR)
of each independent emitter. Firebrands use the structural empirical
flight-distance distribution, Eulerian accumulation, and physical ignition.
Non-structure cells use nonburnable fuel model 93, so vegetative surface
spread is absent and every downwind ignition must be caused by deposited
firebrands.  The global vegetative surface-spread setting switch remains true to suppress vegetative surface-fire propagation.  The current source independently refreshes initialized model-14 building transient HRR under this setting. The configured time-step scaling is
documented by Qin et al.~\cite{QinEtAl2026}:
\[
 \Delta t=\mathrm{CFL}\frac{\Delta}{U_{\rm wind}},
 \qquad \mathrm{CFL}=1,
 \qquad U_{\rm wind}=\SI{17.9}{m.s^{-1}}\;(\SI{40}{mph}).
\]
The fractional step is required to hold this spatial sweep at fixed wind Courant--Friedrichs--Lewy (CFL).
Preprocessing sets each generated stop to
\(\lceil150000/\Delta t\rceil\Delta t\), then records the requested stop,
aligned stop, and integer step count in the simulation specification.

For each represented downwind structure column \(s\), postprocessing takes the
median positive finite arrival time \(t_s\). At least three ignited columns are
required. Mean community ROS is the least-squares slope
\[
 ROS_\Delta =
 \frac{\sum_s(t_s-\bar t)(x_s-\bar x)}
      {\sum_s(t_s-\bar t)^2}.
\]
Because the 10-m grid exactly resolves both physical length scales, its fitted
ELMFIRE result is the comparison anchor,
\[
 ROS_{10}=ROS_{\Delta=\SI{10}{m}}.
\]
The dashed reference line and the 10-m ELMFIRE point therefore use the same
calculated value; no rounded external estimate is substituted.

\section{Derivation of expected behavior}
A resolution-independent structure treatment would conserve each house's
physical footprint, integrated heat release, firebrand source, accumulated
load, and ignition state. Subject to numerical error, it would give
\(ROS_\Delta\approx ROS_{10}\) for every resolution and enough downstream
ignitions to fit all 15 ROS values.

The current cell-local representation does not conserve those quantities. Fine
meshes fragment houses into multiple emitters, intermediate meshes alter
footprints and gaps, and some coarse meshes merge structures or remove
firebreaks. The expected diagnostic result is therefore a nonuniform ROS curve,
including possible failure to propagate at some resolutions. That observation
exposes the defect; against the mesh-independence requirement, its correct
overall status is FAIL.

% BEGIN EXPLICIT EXPECTED-RESULT DERIVATION
The period of the one-dimensional physical indicator is
\(a+d=10+10=20~\mathrm{m}\).  The key raster substitutions are therefore:
\[
\begin{array}{c|c|c}
\Delta~(\mathrm{m})&\text{mode-resampled consequence}&
\text{component-area reference}\\\hline
7&\text{one or two selected cells per coordinate}&
7^2=49\ \text{to}\ (2\cdot7)^2=196~\mathrm{m^2}\\
10&\text{one structure cell followed by one gap cell}&
10^2=100~\mathrm{m^2}\\
30&\text{one selected cell per represented structure}&
30^2=900~\mathrm{m^2}
\end{array}
\]
At \(\Delta=20\) m, each cell contains 10 m of structure and 10 m of gap, so
\(f_i=0.5\) and the \(f_i\geq0.5\) rule selects every cell; the represented
firebreak vanishes.  The same 50-percent selection occurs for the 40-m map.
These are reproducible consequences of the generated mode rule, not inferred
from the output fire spread.

With \(U=40(0.44704)=17.8816~\mathrm{m\,s^{-1}}\), each generated step is
\[
 \Delta t=\Delta/U;
 \quad \Delta t_7=0.391463,\quad
 \Delta t_{10}=0.559234,\quad
 \Delta t_{30}=1.677702~\mathrm{s}.
\]
A truly resolution-independent treatment would give all 15 finite fits,
\(f_{\mathrm{fit}}=15/15=1\),
\(ROS_\Delta=ROS_{10}\), \(CV_R=0\), and \(D_{\max}=0\).
The current native cell treatment is expected instead to produce missing or
nonuniform fits because source HRR scales as \(\Delta^2\) and the represented
geometry changes.  The predeclared mesh-independence requirements permit
\(CV_R\leq0.25\) and \(D_{\max}\leq0.50\); exceeding either limit or failing
any of the 15 fits yields the expected diagnostic FAIL.

The primary publication provides the independent propagation reference calculation used only
for strict stalled-final compatibility: Figure~9 reports zero ROS at 1, 3, 7,
13, 17, 23, 25, and 27~m, and finite propagation at 5, 10, 15, 20, 30, 40,
and 50~m~\cite{QinEtAl2026}.  An early stalled-final dump may therefore use
the no-propagation exception only for the first set, and only after the
meteorology-interval reconstruction and timestep-grid checks also pass.  This
reference calculation does not replace the current-run metrics or authorize accepting an
arbitrary overrun timestamp.
% END EXPLICIT EXPECTED-RESULT DERIVATION

\section{Verification metrics and acceptance criteria}

Preprocessing passes only when all 15 resolutions are generated, the exact
10-m map has no footprint or gap error, the 7-m footprint range is
49--\SI{196}{m^2}, and the median 30-m footprint is \SI{900}{m^2}.
These checks establish that execution comparison uses the intended
resolution-dependent rasterizations.

A complete execution decision additionally requires every input-matched
simulation condition and a finite 10-m ROS anchor. Let \(N_{\rm fit}\) be the number of
resolutions with at least three ignited downwind structure columns. ROS
extraction completeness is
\[
 f_{\rm fit}=\frac{N_{\rm fit}}{15},
\]
and must equal one. For the finite fitted values, mesh dependence is measured
by the coefficient of variation
\[
 CV_R=\frac{\operatorname{std}_{\Delta}(ROS_\Delta)}
            {\operatorname{mean}_{\Delta}(ROS_\Delta)}
\]
and by the largest deviation from the exact-geometry anchor,
\[
 D_{\max}=\max_\Delta\frac{|ROS_\Delta-ROS_{10}|}{ROS_{10}}.
\]
The permissive engineering limits are \(CV_R\le0.25\) and
\(D_{\max}\le0.50\). Missing fits are handled by \(f_{\rm fit}\); they are not
silently discarded and counted as convergence. The overall decision is the
conjunction of every geometry, execution, anchor, fit-completeness, and
mesh-independence check.

\subsection{Justification of metric quantification}

The geometry limits are exact consequences of rasterization.  At
\(\Delta=\SI{10}{m}\), the structure side and gap are both exactly one cell,
so the 1\% footprint and separation allowances cover only file precision and
connected-component bookkeeping.  At \SI{7}{m}, a represented component can
contain one to four cells, giving areas from \(7^2=49\) to
\(4\times7^2=196\ \mathrm{m^2}\).  At \SI{30}{m}, a one-cell represented
structure has area \(30^2=900\ \mathrm{m^2}\).  These checks establish that
the intended mesh-dependent maps were generated; they are not claims of
geometric accuracy.

For the behavioral requirement, \(f_{\rm fit}=1\) is non-negotiable because
mesh independence cannot be inferred after silently discarding resolutions
that fail to propagate.  The \(CV_R\le0.25\) and \(D_{\max}\le0.50\)
limits are intentionally permissive engineering bounds, substantially larger
than ordinary front-fit quantization.  Even with those allowances, a method
with factor-of-two variation or missing fits fails, which is the defect this
case is designed to expose~\cite{QinEtAl2026}.  The 10-m reference is the
actual \(ROS_{10}\) extracted by the same postprocessor, so its plotted
reference and ELMFIRE point must coincide by construction; it is not an
independent nominal value.  Current input identity records and all 15 runs are required
to prevent stale or selectively complete outputs from weakening the expected
FAIL diagnosis.

\subsection{Predeclared acceptance criteria}

Table~\ref{tab:acceptance-contract} summarizes the metrics fixed before
inspection of the calculated results. Numerical values and component statuses
are reported separately in Section~7.

\begin{table}[H]
\centering
\footnotesize
\caption{Predeclared verification metrics and acceptance criteria.}
\label{tab:acceptance-contract}
\begin{tabular}{@{}p{0.23\linewidth}p{0.47\linewidth}p{0.21\linewidth}@{}}
\toprule
Metric & Output, region, and aggregation & Acceptance criterion\\
\midrule
Prepared simulation conditions & Generated simulation specification count divided by the 15 requested grid spacings. & 100\% \\
10-m footprint error & Maximum relative error of represented component areas from \SI{100}{m^2}. & \(\le0.01\) \\
10-m separation error & Maximum relative gap-width error from \SI{10}{m}. & \(\le0.01\) \\
7-m component areas & Minimum and maximum represented connected-component areas. & 49--\SI{196}{m^2} \\
30-m footprint error & Relative error of median represented footprint from \SI{900}{m^2}. & \(\le0.01\) \\
Execution completeness & Current-input identity record final time of arrival (TOA) output and successful termination. & 15 of 15 simulation conditions \\
10-m ROS anchor & Structure-column fit from the current 10-m TOA raster. & Finite \(ROS_{10}\) \\
ROS extraction completeness & Resolutions with at least three ignited downwind columns. & \(f_{\rm fit}=1\) \\
ROS mesh-dependence CV & Coefficient of variation over finite fitted ROS values. & \(CV_R\le0.25\) \\
Maximum anchor deviation & Maximum finite-grid relative deviation from \(ROS_{10}\). & \(D_{\max}\le0.50\) \\
Overall decision & Every geometry, completeness, and ROS component is conjunctive. & PASS only if all pass \\
\bottomrule
\end{tabular}
\end{table}

\section{Simulation configuration and predicted outcome}

Figure~\ref{fig:input-configuration} records the prepared spatial inputs for a 10 m grid.
These panels show the prepared spatial fields, remove the
two-cell numerical buffer, and retain map coordinates in metres.  The left panel shows
the most informative fuel or structure layout available for the case; the right panel
shows the cells selected by the non-positive initial level-set field, making a localized
ignition or firebrand-emission source explicit.

\begin{figure}[H]
\centering
\EvidenceFigure[width=0.98\textwidth,height=0.42\textheight,keepaspectratio]{../figures/input_configuration.pdf}
\caption{Whole-domain prepared input configuration for the representative simulation condition.  The
physical domain is shown without the numerical halo; coordinates, configured wind values,
fuel or structure layout, and initial ignition/source geometry are identified in the panels.}
\label{fig:input-configuration}
\end{figure}

\begin{table}[H]
\centering\footnotesize
\caption{Configuration of the two-dimensional resolution sweep.}
\begin{tabular}{@{}p{0.22\linewidth}p{0.27\linewidth}p{0.13\linewidth}p{0.29\linewidth}@{}}
\toprule
Parameter & Configured value & Units & Role\\
\midrule
Physical domain & \(1200\times400\), flat & m & permits many downwind ignitions\\
Structure side / gap & 10 / 10 & m & exact physical scales at the 10-m grid\\
Initial ignition & complete first structure column & -- & uniform upwind source across all structure rows\\
Grid spacings & 1, 3, 5, 7, 10, 13, 15, 17, 20, 23, 25, 27, 30, 40, 50 & m & exposes raster and cell-local dependence\\
Numerical halo & 2 on every side & cells & excluded from physical geometry\\
Wind & 40 from 270 degrees & mph & constant west-to-east transport\\
Time step & \(\Delta/U_{\rm wind}\) & s & holds wind CFL at one\\
Design-fire peak & 400 & \si{kW.m^{-2}} & structural heat-release rate per unit area (HRRPUA) input\\
Generation rate & 10 & \si{pcs.MW^{-1}.s^{-1}} & cell-local building source\\
Transport / accumulation & empirical / Eulerian & -- & current structural spotting path\\
Ignition delays & 42 plus 300 & s & small-flame and development stages\\
Vegetative surface spread & excluded by fuel 93 & -- & isolates firebrand-driven propagation\\
Stop time & at least 150000 & s & rounded up to an exact timestep multiple\\
Final output & one TOA per simulation condition & -- & source for structure-column ROS fits\\
\bottomrule
\end{tabular}
\end{table}

The dashed reference line is predicted to intersect the 10-m ELMFIRE point
exactly because both are \(ROS_{10}\). A mesh-independent implementation would
cluster all other finite ROS points around this line and provide a fit at every
resolution. The current independent-cell treatment is instead expected to
produce missing fits and large variations, causing the behavioral verification
to fail.

\subsection{Preprocessing, execution, and postprocessing}

input preparation procedure places customizable constants at its beginning,
constructs the periodic map analytically, assigns persistent structure IDs,
writes co-registered GeoTIFF inputs and case-specific fuel tables, and records a
SHA-256 input identity record for every simulation condition. It also writes
structure geometry measurements, making the map checks independently
auditable.

All 15 prescribed conditions are simulated independently and sequentially.
Each result must match the corresponding input configuration before the
resolution comparison is evaluated.

result analysis procedure rejects stale or incomplete outputs, reads the
final time-of-arrival raster, groups cells by structure, fits ROS, calculates
all metrics above, and produces the report figure.

\section{Actual simulation results}

Figure~\ref{fig:domain-result} shows the representative time-of-arrival field from a 10 m grid.
The displayed field is taken from the simulated results, masks missing values, and excludes
the two-cell numerical buffer.  The 10-m field is the reference-scale spatial result used to interpret mesh-dependent structure spread across the resolution sweep.

\begin{figure}[H]
\centering
\EvidenceFigure[width=0.96\textwidth,height=0.50\textheight,keepaspectratio]{../figures/domain_result.pdf}
\caption{Whole-domain time-of-arrival field from the representative completed ELMFIRE run.  This
domain-scale result supplements, rather than replaces, the higher-level verification
profiles, convergence plots, ensemble statistics, and calculated acceptance metrics below.}
\label{fig:domain-result}
\end{figure}

All 15 input-matched ELMFIRE processes reached the source success
message.  The strict terminal-evidence selector accepts
\CompletedVariantCount{} of \RequestedVariantCount{} outputs; process exit alone is
not used as evidence.  Regular terminal dumps are accepted for 1, 3, 5, 7, 10,
20, 30, and 40~m, and the publication-derived independent prediction of no propagation permits the
timestep-aligned stalled-final path for 13, 17, and 25~m.  The selected files,
configured controls, final times, reconstructed pre-jump times, grid residuals,
and rejection reasons are recorded per simulation condition in
calculated verification and validation metrics. The overall decision calculated by postprocessing is
\textbf{\OverallStatus}.

The upper row of Figure~\ref{fig:spatial-resolution} shows the actual
preprocessed geometries. Its lower panel contains fitted ELMFIRE ROS values.
Red crosses at zero identify completed simulations with fewer than three
ignited downwind columns, for which ROS is deliberately not fitted. The dashed
line is \(ROS_{10}=\DxTenReferenceROS~\si{m.s^{-1}}\), exactly equal to the
10-m ELMFIRE point.

\begin{figure}[H]
\centering
\EvidenceFigure[width=\linewidth]{spatial_resolution_sweep.pdf}
\caption{Two-dimensional spatial-resolution diagnostic. Top: mode-resampled
WUI geometry at 7, 10, and \SI{30}{m}. Bottom: mean community ROS from current
ELMFIRE TOA outputs; red crosses identify resolutions without enough downwind
ignitions for a fit, and the dashed line is the run-derived 10-m anchor. The
geometry selection follows the documented comparison
configuration~\cite{QinEtAl2026}.}
\label{fig:spatial-resolution}
\end{figure}

\section{Calculated metrics and verification decision}

\begin{table}[H]
\centering\small
\caption{Calculated verification metrics and component decisions.}
\begin{tabular}{@{}p{0.37\linewidth}p{0.18\linewidth}p{0.17\linewidth}p{0.14\linewidth}@{}}
\toprule
Metric & Acceptance limit & Calculated & Status\\
\midrule
\MetricRows
\bottomrule
\end{tabular}
\end{table}

The postprocessor reports \textbf{\OverallStatus}. Geometry preparation passes,
and all 15 processes executed, but strict terminal evidence is complete for only
\CompletedVariantCount{} of \RequestedVariantCount{} simulation conditions.  The 10-m
anchor, ROS-evaluable fraction, coefficient of variation, and maximum anchor
deviation therefore remain \texttt{N/A}; calculating them from a selected subset
would silently change the predeclared verification population.

\section{Technical interpretation}
The accepted subset contains finite fitted values of approximately
\SI{0.00512}{m.s^{-1}} (3~m), \SI{0.00700}{m.s^{-1}} (5~m),
\SI{0.00770}{m.s^{-1}} (7~m), \SI{0.01123}{m.s^{-1}} (10~m),
\SI{0.02617}{m.s^{-1}} (20~m), \SI{0.01612}{m.s^{-1}} (30~m), and
\SI{0.04036}{m.s^{-1}} (40~m).  These values visibly demonstrate mesh
sensitivity, but they are diagnostic observations rather than a complete
acceptance calculation.

Four successful processes do not provide admissible terminal evidence.  For
15~m, the recorded final time is 152833.592~s and subtracting the configured
3600-s meteorology interval gives 149233.592~s; the timestep-grid residual is
0.302~s.  For 23 and 27~m the corresponding reconstructed times are
146507.060 and 148242.909~s, with grid residuals 0.636 and 0.705~s.  Those
residuals exceed the 0.150-s single-precision-scale tolerance.  The 50-m
reconstruction is grid-aligned (0.023-s residual), but occurs at 148236.198~s,
more than one 2.796-s timestep before the configured stop.  The primary
publication reports propagation at 15 and 50~m, so neither early stall may be
reclassified as independently expected no propagation.  Conversely, the
publication reports no propagation at 23 and 27~m, but that reference calculation does not
waive the timestep-grid requirement.

The current source also produces finite spread at 3 and 7~m while the published
Figure~9 places those resolutions at zero, and it stalls at 15 and 50~m where
that figure shows finite spread.  The resulting \textbf{NOT EVALUABLE} status
is therefore an implementation/reference and terminal-time discrepancy, not a
missing execution and not a reason to weaken the evidence selector.  The runs
used executable SHA-256
\texttt{c225b2e0\allowbreak 0349aa98\allowbreak 128285b3\allowbreak d3ff9b20\allowbreak 31d0e209\allowbreak 44cd9f1c\allowbreak 5ee96595\allowbreak d199a920}
built from source commit
\texttt{a2dfbcdf\allowbreak 72209733\allowbreak c000e5d3\allowbreak 431e4472\allowbreak 3e15efea}; the source worktree already
contained unrelated local modifications, so the commit and executable digest
are both reported for reproducibility.

\section{Scope and limitations}
This case verifies numerical mesh dependence for one idealized periodic WUI
array, one wind, and one structural design fire. It does not validate the
physical accuracy of any individual ROS value. The 10-m result is an internal
run-derived anchor, not experimental validation data. A future
resolution-independent structure implementation should be tested with the same
geometry and acceptance criteria; it should provide all 15 ROS fits and reduce
both variation metrics below their limits.

\begin{thebibliography}{9}
\bibitem{Qin2025}
Y. Qin, \emph{A Physics-Based Eulerian Framework for Modeling Firebrand
Showering in Regional-Scale Wildland and Wildland--Urban Interface Fire
Simulations}, Ph.D. dissertation, University of Maryland, College Park, 2025.

\bibitem{QinEtAl2026}
Y. Qin et al., ``Simulations of firebrand-driven fire spread in
landscape-scale Wildland--Urban-Interface and urban conflagration models,''
\emph{Fire Safety Journal}, vol.~162, 104686, 2026,
doi:10.1016/j.firesaf.2026.104686.
\end{thebibliography}
